\section{Функции приложения}

\subsection{Общие правила}

\begin{enumerate}
  \item Сотрудник видит текущую операцию, товар, ячейку и следующее действие.
  \item Сканирование само по себе не меняет данные: сотрудник подтверждает результат отдельной кнопкой.
  \item Ошибка объясняется простым текстом и не приводит к изменению остатка или результата операции.
  \item После успешного действия сотрудник видит понятный итог и может перейти к следующему заданию.
\end{enumerate}

\subsection{Рабочие сценарии}

Каждый сценарий выполняется по локально загруженному заданию. Итог подтверждается сотрудником и сохраняется на устройстве.

\begin{tabularx}{\textwidth}{L{36mm}Y}
\toprule
\textbf{Операция} & \textbf{Подробный сценарий} \\
\midrule
Приемка & Сотрудник принимает поставку от курьера: сканирует полученные коробки, сверяет их с накладной, подтверждает приемку и при необходимости фотографирует подписанную накладную в приложении. После завершения приемки коробки становятся доступны для операции размещения во внутреннем складе. Подробнее в \hyperref[scenario:receiving]{разделе~\ref{scenario:receiving} «Приемка»}. \\
Размещение & Сотрудник получает задание на размещение товаров. Для каждого товара он сканирует штрихкод товара, затем сканирует ячейку, в которую кладёт товар. После успешного сканирования ячейки товар считается размещённым на полке; сотрудник повторяет действия для следующего товара, пока не будут размещены все товары задания. Подробнее в \hyperref[scenario:placement]{разделе~\ref{scenario:placement} «Размещение»}. \\
\bottomrule
\end{tabularx}

\subsection{Условия ошибок}

\begin{tabularx}{\textwidth}{L{50mm}Y}
\toprule
\textbf{Ситуация} & \textbf{Что должен увидеть и сделать пользователь} \\
\midrule
Код не распознан & Сообщение «Не удалось считать штрихкод. Наведите камеру еще раз»; сотрудник повторяет сканирование. \\
Товар не найден & Сообщение «Товар не найден в задании»; сотрудник повторяет сканирование или возвращается к заданию без сохранения изменений. \\
Товар не подходит к заданию & Сообщение о несоответствии товару или операции; сотрудник не может подтвердить действие до сканирования корректного товара. \\
Нет доступа к камере & Объяснение необходимости доступа и кнопки «Повторить запрос» и «Открыть настройки». \\
Пользователь отменил операцию & Приложение возвращает сотрудника к заданию; локальные данные операции не изменяются. \\
Коробка уже обработана & Сообщение «Коробка уже обработана»; количество, остатки и отчёт о приёмке не меняются. \\
\bottomrule
\end{tabularx}

\subsection{Отчёт о приёмке}

После закрытия приёмки приложение сохраняет на устройстве отчёт о приёмке. В нём указываются номер задания, дата и время, список обработанных коробок, а также все отклонения. Для каждой проблемной ситуации --- например, неизвестного штрихкода, отсутствующей, лишней или повторно отсканированной коробки --- сотрудник добавляет комментарий с причиной.

Фотография подписанной накладной прикрепляется к отчёту о приёмке и сохраняется вместе с ним. После сохранения фотография доступна при просмотре отчёта.
